\documentclass[10pt,a4paper]{amsart}
\usepackage{amsmath,amssymb,amsthm,graphicx}
% Retained minimal preamble for the Stacks Project chapter varieties.tex.
% Verbatim excerpt of the chapter preamble: Xy-pic setup (native preamble lines 14-18), the native theorem
% declarations with the mathematical macros (native preamble lines 154-217), and the project own no-op
% definitions of the reference/slogan/history display environments (native wrapper.tex lines 1-3) that the
% counted statement uses for its slogan. Copied mechanically; nothing added, removed or redefined.
% For commutative diagrams we use Xy-pic
\usepackage[all]{xy}

% We use 2cell for 2-commutative diagrams.
\xyoption{2cell}
% Theorem environments.
%
\theoremstyle{plain}
\newtheorem{theorem}[subsection]{Theorem}
\newtheorem{proposition}[subsection]{Proposition}
\newtheorem{lemma}[subsection]{Lemma}

\theoremstyle{definition}
\newtheorem{definition}[subsection]{Definition}
\newtheorem{example}[subsection]{Example}
\newtheorem{exercise}[subsection]{Exercise}
\newtheorem{situation}[subsection]{Situation}

\theoremstyle{remark}
\newtheorem{remark}[subsection]{Remark}
\newtheorem{remarks}[subsection]{Remarks}

\numberwithin{equation}{subsection}

% Macros
%
\def\lim{\mathop{\mathrm{lim}}\nolimits}
\def\colim{\mathop{\mathrm{colim}}\nolimits}
\def\Spec{\mathop{\mathrm{Spec}}}
\def\Hom{\mathop{\mathrm{Hom}}\nolimits}
\def\Ext{\mathop{\mathrm{Ext}}\nolimits}
\def\SheafHom{\mathop{\mathcal{H}\!\mathit{om}}\nolimits}
\def\SheafExt{\mathop{\mathcal{E}\!\mathit{xt}}\nolimits}
\def\Sch{\mathit{Sch}}
\def\Mor{\mathop{\mathrm{Mor}}\nolimits}
\def\Ob{\mathop{\mathrm{Ob}}\nolimits}
\def\Sh{\mathop{\mathit{Sh}}\nolimits}
\def\NL{\mathop{N\!L}\nolimits}
\def\CH{\mathop{\mathrm{CH}}\nolimits}
\def\proetale{{pro\text{-}\acute{e}tale}}
\def\etale{{\acute{e}tale}}
\def\QCoh{\mathit{QCoh}}
\def\Ker{\mathop{\mathrm{Ker}}}
\def\Im{\mathop{\mathrm{Im}}}
\def\Coker{\mathop{\mathrm{Coker}}}
\def\Coim{\mathop{\mathrm{Coim}}}

% Boxtimes
%
\DeclareMathSymbol{\boxtimes}{\mathbin}{AMSa}{"02}

%
% Macros for moduli stacks/spaces
%
\def\QCohstack{\mathcal{QC}\!\mathit{oh}}
\def\Cohstack{\mathcal{C}\!\mathit{oh}}
\def\Spacesstack{\mathcal{S}\!\mathit{paces}}
\def\Quotfunctor{\mathrm{Quot}}
\def\Hilbfunctor{\mathrm{Hilb}}
\def\Curvesstack{\mathcal{C}\!\mathit{urves}}
\def\Polarizedstack{\mathcal{P}\!\mathit{olarized}}
\def\Complexesstack{\mathcal{C}\!\mathit{omplexes}}
% \Pic is the operator that assigns to X its picard group, usage \Pic(X)
% \Picardstack_{X/B} denotes the Picard stack of X over B
% \Picardfunctor_{X/B} denotes the Picard functor of X over B
\def\Pic{\mathop{\mathrm{Pic}}\nolimits}
\def\Picardstack{\mathcal{P}\!\mathit{ic}}
\def\Picardfunctor{\mathrm{Pic}}
\def\Deformationcategory{\mathcal{D}\!\mathit{ef}}
\newenvironment{reference}{}{}
\newenvironment{slogan}{}{}
\newenvironment{history}{}{}
\usepackage{cleveref}
\crefname{originalreference}{original reference}{original references}
\Crefname{originalreference}{Original reference}{Original references}
\makeatletter
\newlabel{coherent-section-cohomology-projective-space}{{01XS}{source}{}{sourcepointer-1}{}}
\newlabel{coherent-section-cohomology-projective-space@cref}{{[originalreference][0][]01XS}{[1][0][]source}{}{}{}}
\makeatother
\crefname{originalreference}{original reference}{original references}
\Crefname{originalreference}{Original reference}{Original references}
\makeatletter
\newlabel{constructions-lemma-closed-in-projective-space}{{03GL}{source}{}{sourcepointer-2}{}}
\newlabel{constructions-lemma-closed-in-projective-space@cref}{{[originalreference][0][]03GL}{[1][0][]source}{}{}{}}
\makeatother
\crefname{originalreference}{original reference}{original references}
\Crefname{originalreference}{Original reference}{Original references}
\makeatletter
\newlabel{divisors-definition-zero-scheme-s}{{02OQ}{source}{}{sourcepointer-3}{}}
\newlabel{divisors-definition-zero-scheme-s@cref}{{[originalreference][0][]02OQ}{[1][0][]source}{}{}{}}
\makeatother
\crefname{originalreference}{original reference}{original references}
\Crefname{originalreference}{Original reference}{Original references}
\makeatletter
\newlabel{constructions-lemma-surjective-graded-rings-generated-degree-1-map-proj}{{01N1}{source}{}{sourcepointer-4}{}}
\newlabel{constructions-lemma-surjective-graded-rings-generated-degree-1-map-proj@cref}{{[originalreference][0][]01N1}{[1][0][]source}{}{}{}}
\makeatother
\title{Every hyperplane in projective space is isomorphic to projective space of one lower dimension with the twist matched}
\author{The Stacks Project Authors}
\date{stacks-project:varieties.tex:local-snapshot-20261011; GNU FDL 1.2 or later}
\begin{document}
\maketitle
\noindent\textit{Editorial note.} All delivered mathematical text is the author's own, copied verbatim from the pinned Stacks Project chapter file varieties.tex; no statement or proof step was rewritten, reordered or removed.\par\medskip
\noindent\textit{Editorial note.} The counted span cites four cross-chapter labels (including a definition and a section), each recorded in assembly.external\_references with its canonical tag as a navigation pointer.\par\medskip
\noindent\textit{Editorial note.} The proof cites no bibliography item inside the counted span.\par\medskip
\noindent\textit{Allowed standard prerequisites (editorial).}\begin{itemize}
\item Zero scheme of a section and closed subschemes of projective space as Proj of a graded quotient. Stacks Project, Divisors, Definition divisors-definition-zero-scheme-s (tag 02OQ); Constructions, Lemma constructions-lemma-closed-in-projective-space (tag 01M1). Identify the hyperplane with Proj of the quotient of the polynomial ring by the homogeneous ideal of the section.
\item Functoriality of Proj and twists for surjective degree-one generated graded maps. Stacks Project, Constructions, Lemma constructions-lemma-surjective-graded-rings-generated-degree-1-map-proj (tag 01N4). Transfers the graded isomorphism to an isomorphism of Proj with matching twists.
\end{itemize}\medskip
\ifdefined\hypertarget\hypertarget{sourcepointer-1}{}\fi
\noindent\textit{Source reference.} Original reference 01XS: sources/stacks-chow/tags.txt, source lines 2531--2531. Cross-chapter standard prerequisite cited by the author inside the counted span; it is used as a pointer and its own proof is not part of this item.\par\medskip
\ifdefined\hypertarget\hypertarget{sourcepointer-2}{}\fi
\noindent\textit{Source reference.} Original reference 03GL: sources/stacks-chow/tags.txt, source lines 4532--4532. Cross-chapter standard prerequisite cited by the author inside the counted span; it is used as a pointer and its own proof is not part of this item.\par\medskip
\ifdefined\hypertarget\hypertarget{sourcepointer-3}{}\fi
\noindent\textit{Source reference.} Original reference 02OQ: sources/stacks-chow/tags.txt, source lines 3525--3525. Cross-chapter standard prerequisite cited by the author inside the counted span; it is used as a pointer and its own proof is not part of this item.\par\medskip
\ifdefined\hypertarget\hypertarget{sourcepointer-4}{}\fi
\noindent\textit{Source reference.} Original reference 01N1: sources/stacks-chow/tags.txt, source lines 2142--2142. Cross-chapter standard prerequisite cited by the author inside the counted span; it is used as a pointer and its own proof is not part of this item.\par\medskip
\begin{lemma}
\label{lemma-hyperplane}
Let $k$ be a field. Let $n \geq 1$.
Let $i : H \to \mathbf{P}^n_k$ be a hyperplane.
Then there exists an isomorphism
$$
\varphi : \mathbf{P}^{n - 1}_k \longrightarrow H
$$
such that $i^*\mathcal{O}(1)$ pulls back to $\mathcal{O}(1)$.
\end{lemma}

\begin{proof}
We have $\mathbf{P}^n_k = \text{Proj}(k[T_0, \ldots, T_n])$.
The section $s$ corresponds to a homogeneous form in $T_0, \ldots, T_n$
of degree $1$, see
Cohomology of Schemes, Section
\ref{coherent-section-cohomology-projective-space}.
Say $s = \sum a_i T_i$.
Constructions, Lemma \ref{constructions-lemma-closed-in-projective-space}
gives that
$H = \text{Proj}(k[T_0, \ldots, T_n]/I)$ for the graded ideal $I$
defined by setting $I_d$ equal to the kernel of the map
$\Gamma(\mathbf{P}^n_k, \mathcal{O}(d)) \to \Gamma(H, i^*\mathcal{O}(d))$.
By our construction of $Z(s)$ in Divisors, 
Definition \ref{divisors-definition-zero-scheme-s}
we see that on $D_{+}(T_j)$ the ideal of $H$ is generated by
$\sum a_i T_i/T_j$ in the polynomial ring
$k[T_0/T_j, \ldots, T_n/T_j]$. Thus it is clear that $I$ is the ideal
generated by $\sum a_i T_i$. Note that
$$
k[T_0, \ldots, T_n]/I = k[T_0, \ldots, T_n]/(\sum a_i T_i) \cong
k[S_0, \ldots, S_{n - 1}]
$$
as graded rings. For example, if $a_n \not = 0$, then mapping
$S_i$ equal to the class of $T_i$ works. We obtain the desired isomorphism
by functoriality of $\text{Proj}$.
Equality of twists of structure sheaves follows for example from
Constructions, Lemma
\ref{constructions-lemma-surjective-graded-rings-generated-degree-1-map-proj}.
\end{proof}

\end{document}
